% Replacement text for Appendices A and B. The coordinator owns integration.
% Existing equation and section labels are retained where possible.
\section{Observation equivalence and minimum order}\label{app:realization}
A matrix realization $(R,T,B)$ is minimal if its reachable space
$\operatorname{span}\{T^jBv:j\ge0\}$ is the whole state space and its
unobservable space $\{v:RT^jv=0\text{ for all }j\ge0\}$ is zero.
These are linear conditions; they impose no sign restriction on a change
of realization coordinates.

\begin{lemma}[Finite equality certificate and pinned similarity]\label{lem:orbit}
Realizations of dimensions $n,n'$ have the same all-time kernel if and only
if $RT^jB=R'T'^jB'$ for $0\le j<n+n'$. Equivalent minimal realizations have
equal dimensions and a unique similarity intertwining their input,
generator and output maps. For the resolved reverse pair in
Sec.~\ref{sec:model}, after fixing the visible order this similarity has
the form
\begin{equation}
 S=\diag(I_2,U),\qquad U\1=\1,\qquad Q'=S^{-1}QS.\label{eq:orbit}
\end{equation}
Conversely, every such similarity yielding an admissible generator
preserves the complete kernel. The entries of $U$ need not be nonnegative.
\end{lemma}
\begin{proof}
Apply Cayley--Hamilton to $\diag(T,T')$. Vanishing of the first $n+n'$
differences forces every subsequent difference to vanish, and the
convergent exponential series gives all-time equality. Differentiation at
zero proves the converse.

For minimal realizations define
\begin{equation}
 S\left(\sum_jT'^jB'v_j\right)=\sum_jT^jBv_j.
\end{equation}
If the vector on the left is zero, equality of Markov parameters makes
every $RT^\ell$ applied to the vector on the right zero; source
observability makes that vector zero. The same argument in reverse gives
an inverse. Reachability then gives $SB'=B$, $ST'=TS$, $RS=R'$, and
uniqueness. A sufficiently large block Hankel matrix $(RT^{i+j}B)_{i,j}$
has rank equal to the minimal dimension. Every realization of the same
kernel factors this matrix through its state space, so a competing
dimension cannot be smaller.

The zero-time data determine $q_{xy}=\psi_{-,+}(0)$ and
$q_{yx}=\psi_{+,-}(0)$. After putting the same two visible states first,
the reset maps coincide and the columns of $B=B'$ are positive multiples
of the two visible coordinate vectors. Thus $SB=B$ pins the visible
columns of $S$, and $RS=R$ pins its visible rows. This proves the block
form of $S$.

The killed generator $T$ is invertible. Indeed, in a finite irreducible
chain each initial state has a positive probability of producing an
observed jump in some common finite time. Taking the minimum over the
finitely many states gives a strictly positive probability. The Markov
property then bounds survival geometrically over successive intervals,
so $e^{Tt}\to0$ and $T$ is Hurwitz. Since $T\1=-B\1$,
\begin{equation}
 TS\1=ST'\1=-SB\1=-B\1=T\1,
\end{equation}
and invertibility gives $S\1=\1$. The observed jump matrices have nonzero
entries only in the fixed visible block, so they commute with $S$.
Consequently the similarity for $T$ is also the stated similarity for
$Q$. Direct substitution in $Re^{Tt}B$ proves sufficiency.
\end{proof}

For the balanced source the hidden blocks are
\begin{equation}
 X=\begin{pmatrix}2&0&8\\0&7&11\end{pmatrix},\quad
 Y=\begin{pmatrix}3&0\\0&6\\z&z\end{pmatrix},\quad
 H=\begin{pmatrix}-7&4&0\\5&-11&0\\0&0&-2z\end{pmatrix}.
 \label{eq:blocks}
\end{equation}
Writing $Y_j$ for its columns, direct determinants give
\begin{align}
 \det(Y_1,Y_2,HY_1)&=-9z(4z-9),\\
 \det(Y_1,Y_2,HY_2)&=-18z(2z-3),\\
 \det\begin{pmatrix}X_{1,:}\\X_{2,:}\\(XH)_{2,:}\end{pmatrix}
 &=-14(22z+19).
\end{align}
For $z>0$ the first two cannot vanish simultaneously and the last never
vanishes. Hence $(H,Y)$ is reachable and $(X,H)$ observable, both with
hidden dimension three.

To pass from hidden to full minimality, write
$T=\left(\begin{smallmatrix}A_{\rm kill}&X\\Y&H\end{smallmatrix}\right)$.
The columns of $B$ span the visible plane. Applying $T$ to that plane and
subtracting the visible component puts $(0,Y_j)^T$ in the full reachable
space. If $(0,h)^T$ is in that space, applying $T$ and subtracting its
visible component puts $(0,Hh)^T$ there as well. Hidden reachability thus
gives full reachability. If a full vector is unobservable, $Rv=0$ first
forces its visible coordinates to vanish. Successive output equations
then give $XH^jv_H=0$ for every $j$, forcing $v_H=0$ by hidden
observability. Therefore the full minimal order is five for every $z>0$,
including repeated hidden poles and $z=2$. The cap makes every competitor
exactly five dimensional and minimal.

The observed rates and first derivatives determine the visible generator
block $A$, including the killed visible diagonals. After the known row
permutation and rate division, the observed Laplace matrix determines
the visible principal block $K(\lambda)$ of $(\lambda I-T)^{-1}$. For
$\lambda>0$, both $\lambda I-H$ and $\lambda I-T$ are invertible, and
the Schur complement gives
\begin{equation}
 X(\lambda I-H)^{-1}Y
 =\lambda I-A_{\rm kill}-K(\lambda)^{-1}.\label{eq:htransfer}
\end{equation}
Thus every competitor has the same matrix-valued hidden transfer.
Fitting its scalar entries separately would discard this common
realization constraint.

\section{Global support exhaustion and the invariant triangle}\label{app:cone}
\subsection{Modal coordinates and every disconnected hidden support}
Put $\rho=\sqrt6$, $\alpha=9-2\rho$, $\beta=9+2\rho$,
$r_y=4(\rho-1)/5$, and
\begin{gather}
 D=\diag(-\alpha,-\beta,-2z),\quad X_m=(p,q,v),\quad
 Y_m=\begin{pmatrix}r\\f\\w\end{pmatrix},\label{eq:modal}\\
 p=\binom{3+\rho/2}{35\rho/8},\quad
 q=\binom{3-\rho/2}{-35\rho/8},\quad v=\binom8{11},\\
 r=(1,r_y),\quad f=(1,-4(\rho+1)/5),\quad w=(z,z).
\end{gather}
Direct multiplication gives
\begin{equation}
 X(\lambda I-H)^{-1}Y
 =\frac{pr}{\lambda+\alpha}+\frac{qf}{\lambda+\beta}
  +\frac{vw}{\lambda+2z}.\label{eq:residues}
\end{equation}
Here $pr$ and $vw$ are strictly positive rank-one residues, whereas $qf$
has negative off-diagonal entries. In this appendix $z\ge21/2$, so all
three poles are distinct and $2z>\beta>\alpha>0$.

Reciprocity makes hidden connected components the irreducible blocks of
$H'$. On three hidden states there are precisely three component-size
possibilities: $(1,1,1)$, $(2,1)$, and $(3)$. In a disconnected case,
each eigenvalue belongs to exactly one component, since the spectrum is
simple. Every mode has its displayed nonzero transfer residue, by
minimality. A singleton contributes a nonnegative outer-product residue.
Three singletons, or a singleton at $-\beta$, are consequently impossible.
If the singleton were at $-\alpha$, the remaining irreducible pair would
have dominant eigenvalue $-\beta$. Its transfer residue there is
$(X'r_P)(\ell_P^TY')/(\ell_P^Tr_P)\ge0$ by its positive Perron vectors,
again contradicting $qf$. Thus the singleton must have eigenvalue $-2z$.

The remaining pair has zeroth transfer moment
$pr+qf=\diag(6,42)$. For either of its states, reciprocal visible support
makes a positive off-diagonal contribution equivalent to attachment to
both visible states. Nonnegativity therefore forbids a shared visible
neighbor. The two positive diagonal moments require one state attached
only to $x$ and the other only to $y$; neither can lack visible
attachment. Denote their entrance rates by $x_1,x_2$, exits by $y_1,y_2$,
and internal rates by $h_0,k_0$. The zeroth and first moments give
\begin{gather}
 x_1y_1=6,\quad x_2y_2=42,\quad y_1+h_0=7,\quad y_2+k_0=11,\\
 x_1h_0y_2=48,\qquad x_2k_0y_1=105.
\end{gather}
All denominators in the following elimination are positive physical
rates. Multiplying the off-diagonal equations and dividing by the two
diagonal equations gives $h_0k_0=20$. Dividing the first off-diagonal
equation by $x_1y_1=6$ gives
$h_0(11-k_0)=8(7-h_0)$, or $19h_0-h_0k_0=56$. Hence
\begin{equation}
 (h_0,k_0,y_1,y_2,x_1,x_2)=(4,5,3,6,2,7).
\end{equation}
The fast singleton's residue fixes its entrance/exit product; its exit
sum $2z$ removes the remaining scalar freedom and gives entrances
$(8,11)^T$ and exits $(z,z)$. This exhausts all disconnected hidden
supports and leaves exactly the source, up to hidden permutations.

\subsection{Connected supports produce finite nondegenerate triangles}
For a connected hidden support, the hidden block similarity provides a
real invertible $V$ such that
\begin{equation}
 H'=V^{-1}DV\text{ is irreducible Metzler},\qquad
 X'=X_mV\ge0,\qquad Y'=V^{-1}Y_m\ge0.
\end{equation}
The Perron right eigenvector $b=V^{-1}e_1$ has a common strict sign.
Since $X'b=p>0$, that sign is positive. The left Perron eigenvector
$a_P^T=e_1^TV$ also has a common strict sign; $a_P^Tb=1$ makes it
positive. Write $s_i=(a_P)_i>0$ for the first coordinate of column $i$ of
$V$. Sectioning the invariant cone $V\mathbb R_+^3$ at first coordinate
one therefore gives exactly the convex hull of the three columns
$V_{:,i}/s_i=(1,u_i,v_i)^T$. Their affine independence follows from
$\det V\ne0$. Moreover
\begin{equation}
 e_1=\sum_i b_is_i\frac{V_{:,i}}{s_i},\qquad
 b_is_i>0,\quad\sum_i b_is_i=1,
\end{equation}
so the planar origin lies strictly inside this finite triangle. This
argument covers arbitrarily ill-conditioned but invertible similarities;
it assumes no common compactness bound. A zero-first-coordinate ray is
impossible in this connected branch, and disconnected supports have
already been handled separately.

Since $e^{Dt}V=Ve^{H't}$ and $e^{H't}\ge0$, the cone is invariant.
Sectioning after division by the positive first-coordinate contraction,
and rescaling time by $\beta-\alpha$, gives the planar flow
\begin{equation}
 (u,v)\longmapsto(e^{-t}u,e^{-kt}v),\qquad
 k=\frac{2z-9+2\rho}{4\rho}>1.\label{eq:flow}
\end{equation}
Nonnegative input barycentric coordinates and output values mean the
triangle contains $(1,A),(-m,B_0)$ and is contained in the wedge
\begin{gather}
 A=z,\quad B_0=z/r_y,\quad m=(7+2\rho)/5,\\
 L=(3-\rho/2)/8,\quad R_w=35\rho/88,\\
 v\ge L(-m-u),\qquad v\ge R_w(u-1).\label{eq:wedge}
\end{gather}

\subsection{Vertex signs, boundary cases and the reflected inequality}
A wedge point with nonpositive height has abscissa in $[-m,1]$.
If the triangle had only one positive-height vertex with abscissa $u_0$,
positive height at input abscissa $1$ would require $u_0\ge1$: otherwise
every convex combination using that vertex with positive weight has
abscissa less than one. The other positive input similarly requires
$u_0\le-m$, a contradiction. The interior origin requires a
negative-height vertex. Thus the three vertex heights consist of two
strictly positive values and one strictly negative value; zero-height
vertices are excluded. The negative vertex is strictly inside the
abscissa interval, and the two positive vertices must extend to either
end. Consequently they can be labeled
\begin{equation}
 (-\ell,P),\quad(c,Q_v),\quad(d,-W),\quad
 \ell\ge m,\ c\ge1,\ -m<d<1,\ P,Q_v,W>0.\label{eq:vertices}
\end{equation}
At this stage $\ell=m$, $c=1$, and $d=0$ have not been discarded.
Put $a=k-1>0$. Wedge membership gives
\begin{align}
 P&\ge L(\ell-m),\qquad Q_v\ge R_w(c-1),\nonumber\\
 W&\le\min\{L(m+d),R_w(1-d)\}<1.\label{eq:ineq1}
\end{align}
The last strict inequality follows from
$W<L(m+1)<(1/4)(7/2)=7/8$. Lower-face velocities at the two upper
vertices give
\begin{equation}
 (a\ell+kd)P\le\ell W,\qquad
 (ac-kd)Q_v\le cW.\label{eq:ineq2}
\end{equation}
Input containment above and below the relevant edges gives
\begin{align}
 B_0&\le H_c:=\frac{(c+m)P+(\ell-m)Q_v}{\ell+c},\label{eq:height}\\
 A&\le\frac{(c-1)P+(\ell+1)Q_v}{\ell+c},\label{eq:upperA}\\
 (Q_v+W)(1-d)&\le(A+W)(c-d),\label{eq:lowerA}\\
 (P+W)(d+m)&\le(B_0+W)(\ell+d).\label{eq:lowerB}
\end{align}
All geometric denominators here are strictly positive.

If $d<0$, the second inequality~\eqref{eq:ineq2} gives $Q_v<W/a$.
Using explicitly the right-output inequality in~\eqref{eq:ineq1} and
then~\eqref{eq:lowerB}, one obtains
\begin{align}
 A&\le Q_v+\frac{(c-1)P}{\ell+c}
 \le Q_v\left(1+\frac{P}{R_w(\ell+c)}\right)\nonumber\\
 &<\frac{W}{a}\left(1+\frac{B_0+W}{R_w(m+d)}\right)
 \le\frac{L}{a}\left(m+d+\frac{B_0+W}{R_w}\right)\nonumber\\
 &<\frac{L}{a}\left(m+\frac{B_0+1}{R_w}\right)
 <\frac{25/8+(25/18)(z+1)}{z-7}.
\end{align}
In the middle step, $(P+W)/(\ell+c)<(B_0+W)/(m+d)$ follows from
$\ell+d<\ell+c$. The final comparison uses
$m<5/2$, $L<1/4$, $R_w>9/10$, $B_0<z$ and $a>(z-7)/5$.
The last fraction is below $z=A$ for $z\ge10$, since clearing the
positive denominator gives
\begin{equation}
 72(z-10)^2+836(z-10)+835>0.
\end{equation}
This contradiction excludes $d<0$.

For $d\ge0$, $a>7/10$ and~\eqref{eq:ineq2} imply
$P\le W/a<10/7<2<B_0<A$. Upper containment gives $Q_v\ge A$ and
$\ell>m$: otherwise $B_0\le H_c=P$. Since $aQ_v-W>0$, the equivalent
inequality $c(aQ_v-W)\le kdQ_v$ also excludes $d=0$. Write
$\epsilon=1-d>0$. From $c\ge1$ and the same inequality,
\begin{equation}
 \epsilon\le\frac{A+W}{kA}
 \le\frac{A+R_w\epsilon}{kA},\qquad
 \epsilon\le\frac1{k-R_w/A}.\label{eq:apexbound}
\end{equation}
The denominator is positive. The wedge apex satisfies
\begin{equation}
 d_0=\frac{R_w-Lm}{L+R_w}>0,\qquad
 1-d_0=\frac{L(m+1)}{L+R_w}.
\end{equation}
At $z=21/2$ the latter expression exceeds the last bound
in~\eqref{eq:apexbound} by more than $0.014259$, as certified by rational
evaluation in $\mathbb Q(\sqrt6)$. Since $k-R_w/z$ increases with $z$,
every candidate has $d>d_0$. The tighter bottom bound is therefore
$W\le R_w\epsilon$, with equality allowed.

\subsection{The global envelope, including both sides of its pole}
Fix $d,W$ from a feasible triangle. The left bounds imply
\begin{equation}
 L(\ell-m)\le g(\ell):=\frac{\ell W}{a\ell+kd}.
\end{equation}
Here $g'>0$, $g''<0$, and the left-intersection equation for
$\ell_*=m+t$ is a quadratic in $t$ with positive leading coefficient and
negative constant. It has precisely one positive root. Thus
\begin{equation}
 \ell\le\ell_*,\quad P\le P_*:=Lt=g(m+t),\qquad
 Lt[a(m+t)+kd]=(m+t)W.\label{eq:leftmax}
\end{equation}
In particular $P_*<W/a<2$. The partial derivatives
\begin{equation}
 \frac{\partial H_c}{\partial P}=\frac{c+m}{\ell+c}>0,
 \qquad
 \frac{\partial H_c}{\partial\ell}
 =\frac{(c+m)(Q_v-P)}{(\ell+c)^2}>0
\end{equation}
remain strict when $P$ is increased to $P_*$ and $\ell$ to $\ell_*$,
because $Q_v\ge A>2$. This replacement is an upper-bound relaxation;
it need not retain every unused constraint. Equality in this relaxation
forces both $P=P_*$ and $\ell=\ell_*$.

For the right vertex set
\begin{equation}
 \mathcal L(c)=\frac{A+W}{\epsilon}(c-d)-W,
 \qquad \mathcal H(c)=\frac{cW}{ac-kd}.
\end{equation}
The line bounds $Q_v$ for every $c\ge1$. When $c\le kd/a$, the
invariance inequality has nonpositive left coefficient and imposes no
additional bound; one must not divide by it. When $c>kd/a$, it gives
$Q_v\le\mathcal H(c)$. The relevant crossover is
\begin{equation}
 c_* =\frac{k}{a}\left[d+\frac{W\epsilon}{A+W}\right],\qquad
 Q_*^v=\frac{Ad+W}{a\epsilon}.\label{eq:rightmax}
\end{equation}
Indeed $c_*-kd/a=kW\epsilon/[a(A+W)]>0$, and on the positive-denominator
branch the exact factorization is
\begin{equation}
 \mathcal L(c)-\mathcal H(c)
 =\frac{a(A+W)(c-d)(c-c_*)}{\epsilon(ac-kd)}.
\end{equation}
The other algebraic intersection $c=d$ lies outside $c\ge1$.
Therefore the active bound is the line up to $c_*$ and the hyperbola
afterward, including the possible interval below the pole. If $c_*<1$,
the hyperbola at $c=1$ is less than $\mathcal L(1)=A$, and decreases for
larger $c$, contradicting $Q_v\ge A$. Hence $c_*\ge1$ and $Q_*^v\ge A$.

Write $t=\ell_*-m>0$. At fixed left maximum,
\begin{equation}
 H_c=P_*+\frac{t(Q_v-P_*)}{\ell_*+c},\qquad
 \partial_{Q_v}H_c=\frac{t}{\ell_*+c}>0.
\end{equation}
Along the line its derivative with respect to $c$ is
\begin{equation}
 \frac{t[(A+W)(\ell_*+d)/\epsilon+W+P_*]}{(\ell_*+c)^2}>0.
\end{equation}
Along the hyperbola,
\begin{equation}
 \frac{dH_c}{dc}
 =\frac{t(P_*-\mathcal H(c))}{(\ell_*+c)^2}
 -\frac{tWkd}{(\ell_*+c)(ac-kd)^2}<0,
\end{equation}
because $\mathcal H(c)>W/a>P_*$. Thus the unique right maximum is
$(c_*,Q_*^v)$, and equality forces both coordinates to have those values.

It remains to increase $W$ at fixed $d$. Put
$r_c=c_*/Q_*^v=k\epsilon/(A+W)$; the resulting envelope is
\begin{equation}
 H_* =\frac{t[Lm+(1+Lr_c)Q_*^v]}{m+t+r_cQ_*^v}.\label{eq:envelope}
\end{equation}
Implicit differentiation of~\eqref{eq:leftmax} gives $dt/dW>0$, since
the derivative of its left-minus-right expression with respect to $t$,
after using the equality, is
\begin{equation}
 L\left[a(m+t)+\frac{kdm}{m+t}\right]>0.
\end{equation}
Also $dQ_*^v/dW=1/(a\epsilon)>0$ and $dr_c/dW<0$.
For $H_*(t,Q,r)$, the three relevant partial derivatives have numerators
\begin{equation}
 [Lm+(1+Lr)Q](m+rQ),\quad t(m+t+Lrt),\quad tQ(Lt-Q),
\end{equation}
over the common positive square $(m+t+rQ)^2$. Their signs are positive,
positive, and negative. These signs hold throughout increasing $W$ to
$R_w\epsilon$, since $Lt<W/a<2<A\le Q_*^v$. Thus the increase is
strict, and the final bound is attained only at $W=R_w\epsilon$.

\subsection{Discriminant exclusion and propagation of every equality}
After the last saturation, put $J=am+k$ and solve
\eqref{eq:leftmax} for $\epsilon$:
\begin{equation}
 \epsilon=\frac{Lt(J+at)}{R_wm+(kL+R_w)t},\qquad
 H_*-B_0=\frac{f_0+f_1t+f_2t^2}{(J+at)^2}.\label{eq:quadratic}
\end{equation}
All denominators are positive. Direct rational substitution gives
\begin{align}
 f_0&=AR_wm/L-B_0J^2,\\
 f_1&=A(R_w/L+k-am)+LJ^2+R_wJ-2B_0Ja,\\
 f_2&=a[LJ+R_w-A(1+Lk/R_w)-B_0a].\label{eq:fcoeff}
\end{align}
In particular
\begin{align}
 f_0&=-\frac{19+9\rho}{100}z\left(z^2-11z+\frac{211}{44}\right),\label{eq:f0}\\
 f_1^2-4f_0f_2
 &=-\left(\frac{11}{1200}+\frac{11\rho}{4200}\right)
 \left(z+\frac{19}{22}\right)^2\frac{\mathcal P(z)}{352}.
 \label{eq:discriminant}
\end{align}
Descartes' rule and opposite signs at zero and positive infinity imply
that $\mathcal P$ has exactly one positive root. Rational endpoint signs
isolate it in $(10.55714966866,10.55714966867)$. The positive root of the
quadratic in~\eqref{eq:f0} is below this interval: that quadratic is
positive at the lower endpoint and strictly increasing there and
above. Hence for every $z>\zs$, $f_0<0$ and the discriminant is negative.
The quadratic has no real zero and is negative at zero, so it is
negative everywhere. The necessary envelope bound excludes every
connected hidden support, without any boundedness assumption on the
candidate triangle coordinates.

At $z=\zs$, exact interval evaluation gives
\begin{equation}
 -0.520929<f_0<-0.520928,\quad
 5.496508<f_1<5.496509,\quad
 -14.498918<f_2<-14.498917.
\end{equation}
The numerator is nonpositive and vanishes only at
$t_*=-f_1/(2f_2)>0$. Start with any feasible triangle and follow the
three upper-bound relaxations above. Its original upper input height is
at least $B_0$; its final saturated height is at most $B_0$. Therefore
every relaxation is equality. Strict increase under $W$ saturation
first forces the original $W=R_w(1-d)$. The final numerator forces
$t=t_*$. Equation~\eqref{eq:quadratic} now fixes $\epsilon$ and hence
$d,W$ uniquely. Equality in the left and right relaxations fixes
$(\ell,P)=(m+t_*,Lt_*)$ and $(c,Q_v)=(c_*,Q_*^v)$.
This explicitly rules out equality lost in any intermediate relaxation.

The left and bottom vertices lie on the left and right output faces.
Input $(1,A)$ lies on the opposite lower-right edge, and input
$(-m,B_0)$ lies on the opposite upper edge. The paired input/output
equalities remove the same two visible reciprocal pairs. The two
lower-face tangencies remove both directions between the upper
vertices. These are necessary exact zeros; positivity of the remaining
rates is an additional existence check, not a consequence of the zeros
alone.

Finally a fixed triangle leaves no free positive column scaling. For
its sectioned vertex matrix $V$, set
$\bar H=V^{-1}DV$, $\bar Y=V^{-1}Y_m$, $\bar X=X_mV$.
If its columns are rescaled by a positive vector $w$, hidden row
normalization requires $\bar H w+\bar Y\1=0$. Since all three hidden
eigenvalues are nonzero, there is exactly one possible scaling:
\begin{equation}
 w=-\bar H^{-1}\bar Y\1,\quad W_d=\diag(w),\quad
 H'=W_d^{-1}\bar HW_d,\quad X'=\bar XW_d,\quad
 Y'=W_d^{-1}\bar Y.\label{eq:normalize}
\end{equation}
Hidden permutation is the only remaining coordinate freedom. Thus the
equality triangle gives at most one physical connected-hidden generator.
Appendix~\ref{app:certificate} supplies its existence, positive scaling,
exact reciprocal support and entropy separation independently. Together
with the disconnected classification, this proves the endpoint and
supercritical conclusions over the entire stated model class.
